\textbf{结论名称：}直线方程的五种形式（点斜式、斜截式、两点式、截距式、一般式）

\textbf{适用条件：}
点斜式、斜截式要求斜率存在（直线不垂直于\(x\)轴）；两点式要求\(x_1\neq x_2\)且\(y_1\neq y_2\)；截距式要求横、纵截距均存在且不为零；一般式适用于坐标平面内所有直线。

\textbf{变量说明：}
\((x_0,y_0)\)为直线上已知点；\(k\)为斜率；\(b\)为\(y\)轴截距；\((x_1,y_1),(x_2,y_2)\)为直线上两点；\(a,b\)分别为\(x\)轴、\(y\)轴的非零截距；\(A,B,C\)为实数且\(A^2+B^2\neq0\)。

\textbf{核心公式：}
\[
\boxed{
\begin{aligned}
\text{点斜式:}&\quad y-y_0 = k(x-x_0) \\[2pt]
\text{斜截式:}&\quad y = kx + b \\[2pt]
\text{两点式:}&\quad \dfrac{y-y_1}{y_2-y_1} = \dfrac{x-x_1}{x_2-x_1} \\[2pt]
\text{截距式:}&\quad \dfrac{x}{a} + \dfrac{y}{b} = 1 \\[2pt]
\text{一般式:}&\quad Ax + By + C = 0
\end{aligned}
}
\]

\textbf{使用场景：}已知一点与斜率用点斜式；已知斜率与\(y\)截距用斜截式；已知两点用两点式；已知两非零截距用截距式；一般式用于表示任意直线，便于讨论位置关系与距离。

\textbf{简单例子：}已知直线过点\(P(1,2)\)且斜率为\(3\)，按点斜式得\(y-2=3(x-1)\)，化为一般式为\(3x-y-1=0\)。若已知两点\((1,2)\)与\((3,4)\)，由两点式得\(\frac{y-2}{4-2}=\frac{x-1}{3-1}\)，化简得\(x-y+1=0\)。